\documentclass[11pt]{article}
\usepackage[a4paper,margin=22mm]{geometry}
\usepackage{fontspec}
\setmainfont{LinLibertine_R.otf}[BoldFont=LinLibertine_RB.otf,ItalicFont=LinLibertine_RI.otf,BoldItalicFont=LinLibertine_RBI.otf]
\usepackage{xcolor}
\usepackage{hyperref}
\hypersetup{colorlinks=true,urlcolor=blue}
\setlength{\parindent}{0pt}
\setlength{\parskip}{6pt}
\sloppy
\emergencystretch=2em
\begin{document}
\section*{Historical Britannica: extracted source}
Supplement to the third edition of the Encyclopaedia Britannica ... Illustrated with ... copperplates - Volume 1, ABE-IMP - EB.7\par
{\small Source row 1776. Parser: nvidia/nemotron-3.5-parse.\par}
{\small This is model-extracted text, not a corrected transcription. OCR errors, reading order and region alignment require human review.\par}
\hrule\medskip
{\small\itshape \#\# ELECTRICITY.\par}

{\scriptsize\color{gray} Region region-0; provider order 0.\par}
\medskip

{\small\itshape 584\par}

{\scriptsize\color{gray} Region region-1; provider order 1.\par}
\medskip

fed. The experiment fucceeds as well with a metal globe; may, even although the mouth be pretty wide; in which each, there is not a perfect balance of action in every direction. The electrometer may be made to touch the bottom of the globe, or anywhere one too near the mouth, without acquiring any reliable electricity; but if we conduct any feasible with the elettronter, it will infanty be electrical and fromly repelled. Deep cylinders and all round wfefle with narrow months, exhibit the time frameful of electricity within, except near the brims, although frungely electric with: out; and even open metal copa have the interior electronic much diminished.
\par

{\scriptsize\color{gray} Region region-2; provider order 2.\par}
\medskip

Refreshing on this valuable proposition of Mr Caven- dih, we fee clearly why an overcharged electric is only specifically for) and that this will be the case even though we attempt to accumulate a great quantity of electricity in it, by melting it in a thin glass globo, and elechifying its while liquidy and keeping up the acne making force still it because quite cold. The perfect writer, not having considered the fubject with that jurdicions accuracy on Mr Cavendillic exercise, had hopes of producing a powerful and permanent electric in this way, and was mortified and puzzled by the disappointment, till be few his milthake on reading Mr Cavendill’s differentiation.
\par

{\scriptsize\color{gray} Region region-3; provider order 3.\par}
\medskip

These observations alto point out a thing which should be attended to no more experiments for diforerwy the electrically excited in the spontaneous operations of nature, as in chemical composition and decompontition, congellation, future, expansion, etc. It has been usual to put the fullbatters into glubs, or other nonconducting velfle, or into wefle which conduct very imperfectly. In this left cafe especially, the very faint electricity which is probably, infantly forms a compulsion to trifle in the fulchance of the wefl, and the apparatus becomes almost neutral, although there may have been a great deal of electricity excited. It will be proper to consider whether the intent of the expeiriment will admit of metallic reaffle. In the experiments on metalline foliations, the belt method feems to be, to make the veffel itself the fullbearance that is to be diffused.
\par

{\scriptsize\color{gray} Region region-4; provider order 4.\par}
\medskip

For similar reasons we may valedly, without a more minute examination, that bodies of such cases when overcharged, will have the redundant fluid much denser near the furnace than in the interior parts; and denfer in all directions, bumps, projections, angles, and near the exist of oblong bodies, and that is, general, the quantity of redundant fluid, or redundant matter, will be much more nearly proportional to the foreclosed bodies than their quantities of matters. All that is fully proved by experience. The experiment of the electrical chain is a very beautiful one. Lay a long metal chain in an influential metal difarified with with a elettronter. Let one end be held an inch or two above the coil by a file thread. Electricity the whole, and oblove the divergence of the electrometer , thus gradually drawing up the chain from the coil, the electroneter will gradually fall lower, and lowering the chain again will gradually rate to.
\par

{\scriptsize\color{gray} Region region-5; provider order 5.\par}
\medskip

{\small\itshape 10.\par}

{\scriptsize\color{gray} Region region-6; provider order 6.\par}
\medskip

{\small\itshape 10.\par}

{\scriptsize\color{gray} Region region-7; provider order 7.\par}
\medskip

for beginning the computation at the extremity of the prime conductor. It certainly shouldn have been from the extremity. Had the prime conductor been a single globe, it should have begun from the center of this globe. If it was of the usual form, with an out, standing wire, terminated by a large ball, the affection of body of the condition (should certainly have been taken into the account. In short, almost any point of the longeradiduer might have been accommodated to his Lordship's theory.
\par

{\scriptsize\color{gray} Region region-8; provider order 8.\par}
\medskip

We might now proceed to investigate the distribution of the electric fluid in bodies expelled to the action of others, and particularly in the doing conduction made ufc of in our preparatory propositions. The problem is determines, when the helium and diameter of cylinder disc conductors are given; but even when the electric employed for inducing the electricity is in the form of a globe, we must meet any functions of the distances that are pretty complex, and obliges us to have execute to record flux.The most actions of two along conductors, of considerable diameters, give a problem that will occupy the first mathematics; but which is quite interpose for the factory attached. Not is a minor knowledge of the disposition of the fluid of very important service. We may therefore content ourselves where a general representation of the face of the fluid in the following manner, which will give us a pretty di- field notion how it will ask in mok cells.
\par

{\scriptsize\color{gray} Region region-9; provider order 9.\par}
\medskip

Let A (fig. z.o.) be an overcharged sphere, and BC General a conducting cylinder or primary body; draw \_be\_ parched rallet to BC; and credit perpendaciana B, Cc, Pz, dipifice size to represent the capable density of the fluid, were like the conductor is in his natural flame; but let B, Cry, Pz, etc.reprether unequal densities in its different points, while in the vicinity of the overcharged figure. These estimates must be bounded by e line \_in\_, which will cut the line \_b\_ in the point of the perpendicular, drawn from the neutrals point of the conductor. The whole quantity of fluid in the conductor is represented by the parallelogram BC cb; which multt therefore be equal to the flow BC vary; the redanuation fluid in any portion CP or IN is represented by the faces \_crr\_, \_pns\_; and the redundant matter, or deficient fluids in any porous BC, is represented by bdvg. The action of this body on any body placed near it, depends entirely on the area contained between this curveline and its axis \_bc.\_ The only circumstance that we can alteration with respect to this curve is, that the \_variations of curvature\_ in every point are proportional to the forces exerted by the Sphere A; and are therefore inversely as the figures of the distance from A. This property will be demonstrated by and live. The place of n, and the magnitude of the ordinates, will vary as the diameter of the conductor varies. We shall consider this a little more, particularly in four cafes which will occur afterwards. We may consider the impactful that that can occur; namely, when the con ductor is, like a wire, if no flexible diameter, may, as containing only one row of particles.
\par

{\scriptsize\color{gray} Region region-10; provider order 10.\par}
\medskip

Let AE (fig. az.), be facts of derde conducting as in real; and let B, C, E, E, K, exc.reprethe de duality of the fluid would become experts in being kept in this fate of inequable density by the repulsion for fome over charged body. A particle in C is impelled in the dielectric region CE by all the fluid on the tide of A, and in the direction
\par

{\scriptsize\color{gray} Region region-11; provider order 11.\par}
\medskip
\newpage\section*{Source record}
bri-1776-5ed864a29d36\par
Image SHA-256: \texttt{5ed864a29d363db6a7c9380249079c9c}\par\texttt{50b16b15bedde2855c6814920bf98399}\par
Dataset revision: \texttt{309f765b1ecc83e0a7f7}\texttt{762e44b4a927a7555443}\par
National Library of Scotland, Encyclopaedia Britannica Illustrated (CC0).\par
\url{https://huggingface.co/datasets/NationalLibraryOfScotland/encyclopaedia_britannica_illustrated/viewer/default/train?row=1776}
\end{document}
